\section{Verification of \texttt{update} (Tasks 4.1--4.4)}
\label{sec:update}

\subsection{Specification}

Let $a_0$ denote the array state on entry and $a_1$ the array state on exit.

\begin{align*}
  \textbf{Precondition:}\quad &
    \mathit{rev\_sorted}(a_0, n) \;\wedge\;
    \hindex(a_0, n) = h \;\wedge\;
    0 \leq i < n \;\wedge\;
    \forall\, k < n.\; a_0[k] \geq 0 \\
  \textbf{Postcondition:}\quad &
    \mathit{rev\_sorted}(a_1, n) \;\wedge\;
    \hindex(a_1, n) = \mathit{result} \;\wedge\; {} \\
    & \mathit{contents}(a_1) = \mathit{contents}(a_0) - \mset{a_0[i]} + \mset{a_0[i] + 1}
\end{align*}

where $\mathit{contents}(a) = \mset{a[0], a[1], \ldots, a[n{-}1]}$ is the multiset of array values.

\subsection{Preliminary Analysis}

Let $x = a_0[i]$. The function:
\begin{enumerate}
  \item Binary searches $[0, i)$ for the leftmost position $\mathit{lo}$ with $a_0[\mathit{lo}] = x$.
  \item Sets $a[\mathit{lo}] \gets x + 1$.
  \item Returns $h+1$ if $\mathit{lo} = h$ and $x + 1 = h + 1$ (i.e., $x = h$); otherwise returns $h$.
\end{enumerate}

\begin{lemma}[Equal-value block structure]\label{lem:equal-block}
If $\mathit{rev\_sorted}(a, n)$ and $a[i] = x$, then the positions with value $x$ form
a contiguous block $[p, q]$ where $p \leq i \leq q$ and:
\[
  \forall\, j < p.\; a[j] > x \qquad\text{and}\qquad
  \forall\, j > q,\, j < n.\; a[j] < x.
\]
\end{lemma}
\begin{proof}
By reverse-sorting, $a[j] \geq a[i] = x$ for $j \leq i$ and $a[j] \leq a[i] = x$ for $j \geq i$.
Define $p = \min\{j : a[j] = x\}$ and $q = \max\{j : a[j] = x\}$.
For $j < p$: $a[j] \geq x$ (since $j < p \leq i$), and $a[j] \neq x$ (by minimality of $p$),
so $a[j] > x$. Similarly for $j > q$.
Contiguity: for $p \leq j \leq q$, $a[j] \leq a[p] = x$ and $a[j] \geq a[q] = x$, so $a[j] = x$.
\end{proof}

\subsection{Binary Search Invariant}

\begin{invariant}\label{inv:update-search}
At the top of each iteration of the binary search in \texttt{update}:
\[
  K(\mathit{lo}, \mathit{hi}) \;\stackrel{\text{def}}{\iff}\;
  \begin{cases}
    0 \leq \mathit{lo} \leq \mathit{hi} \leq i \\
    \forall\, j.\; 0 \leq j < \mathit{lo} \implies a[j] > x \\
    a[\mathit{hi}] = x
  \end{cases}
\]
where $x = a[i]$.
\end{invariant}

\begin{remark}
The binary search finds the leftmost occurrence of $x$ in $a[0..i]$.
Since $a[i] = x$, the position $\mathit{hi} = i$ initially serves as a witness.
The invariant maintains that some position at or after $\mathit{hi}$ has value $x$,
and all positions before $\mathit{lo}$ have strictly greater values.
\end{remark}

\subsection{Task 4.1: Safety}

\begin{theorem}[Safety of \texttt{update}]\label{thm:update-safety}
Under the precondition, \texttt{update} terminates and performs no out-of-bounds
array accesses.
\end{theorem}

\begin{proof}
\textbf{Initialization of $x$.}
$a[i]$ is accessed with $0 \leq i < n$ (precondition), so within bounds.

\textbf{Binary search array access.}
The access $a[\mathit{mid}]$ requires $0 \leq \mathit{mid} < n$.
From $K(\mathit{lo}, \mathit{hi})$ and $\mathit{lo} < \mathit{hi}$:
$\mathit{mid} = \mathit{lo} + (\mathit{hi} - \mathit{lo})/2$,
so $\mathit{lo} \leq \mathit{mid} < \mathit{hi} \leq i < n$.
Also $\mathit{mid} \geq \mathit{lo} \geq 0$. Thus $0 \leq \mathit{mid} < n$.

\textbf{Post-search assignment.}
$a[\mathit{lo}]$ is accessed after the loop. From $K(\mathit{lo}, \mathit{hi})$ with
$\mathit{lo} = \mathit{hi}$: $0 \leq \mathit{lo} \leq i < n$.

\textbf{Termination.}
Variant $V = \mathit{hi} - \mathit{lo}$. Same argument as Theorem~\ref{thm:opt-safety}:
$V \geq 0$ by invariant, strictly decreases each iteration. At most $\lceil\log_2(i)\rceil$ iterations.

\textbf{Post-search conditional.}
$a[\mathit{lo}]$ in the condition $a[\mathit{lo}] = h + 1$ uses the same index, already shown in bounds.
\end{proof}

\subsection{Binary Search Correctness}

\begin{lemma}[Binary search finds leftmost $x$]\label{lem:search-correct}
After the binary search loop terminates, $\mathit{lo}$ is the smallest index
in $\{0, \ldots, i\}$ with $a[\mathit{lo}] = x$.
\end{lemma}

\begin{proof}
\textbf{Initialization.}
$\mathit{lo} = 0$, $\mathit{hi} = i$.
\begin{itemize}
  \item $0 \leq 0 \leq i$: holds since $i \geq 0$.
  \item $\forall\, j < 0.\; a[j] > x$: vacuously true.
  \item $a[i] = x$: by definition of $x$.
\end{itemize}
So $K(0, i)$ holds.

\textbf{Maintenance.}
Assume $K(\mathit{lo}, \mathit{hi})$ and $\mathit{lo} < \mathit{hi}$.
Let $\mathit{mid} = \mathit{lo} + (\mathit{hi} - \mathit{lo}) / 2$.

\textbf{Case $a[\mathit{mid}] = x$:}
Set $\mathit{hi}' = \mathit{mid}$. Verify $K(\mathit{lo}, \mathit{mid})$:
\begin{itemize}
  \item $0 \leq \mathit{lo} \leq \mathit{mid} \leq i$: since $\mathit{mid} < \mathit{hi} \leq i$.
  \item Left part inherited from $K(\mathit{lo}, \mathit{hi})$.
  \item $a[\mathit{mid}] = x$: by the case condition.
\end{itemize}

\textbf{Case $a[\mathit{mid}] \neq x$:}
Set $\mathit{lo}' = \mathit{mid} + 1$. We must show $a[\mathit{mid}] > x$.
Since $\mathit{mid} \leq i$ and the array is reverse-sorted,
$a[\mathit{mid}] \geq a[i] = x$. Since $a[\mathit{mid}] \neq x$, we get $a[\mathit{mid}] > x$.

Verify $K(\mathit{mid} + 1, \mathit{hi})$:
\begin{itemize}
  \item $0 \leq \mathit{mid} + 1 \leq \mathit{hi} \leq i$: $\mathit{mid} + 1 \leq \mathit{hi}$
    since $\mathit{mid} < \mathit{hi}$.
  \item For $j < \mathit{mid} + 1$:
    If $j < \mathit{lo}$: $a[j] > x$ by $K(\mathit{lo}, \mathit{hi})$.
    If $\mathit{lo} \leq j \leq \mathit{mid}$: by reverse-sorting, $a[j] \geq a[\mathit{mid}] > x$.
  \item $a[\mathit{hi}] = x$: inherited from $K(\mathit{lo}, \mathit{hi})$.
\end{itemize}

\textbf{Postcondition.}
On exit: $\mathit{lo} = \mathit{hi}$. From $K(\mathit{lo}, \mathit{lo})$:
$a[\mathit{lo}] = x$ and $\forall\, j < \mathit{lo}.\; a[j] > x$.
So $\mathit{lo}$ is the leftmost position with value $x$.
\end{proof}

\subsection{Task 4.2: Reverse-Sorted Preservation}

\begin{theorem}[Sorted preservation]\label{thm:update-sorted}
After executing \texttt{update}, $\mathit{rev\_sorted}(a_1, n)$.
\end{theorem}

\begin{proof}
The only modification is $a[\mathit{lo}] \gets x + 1$ where $a_0[\mathit{lo}] = x$.
For all other positions $j \neq \mathit{lo}$: $a_1[j] = a_0[j]$.

We must verify $a_1[j] \geq a_1[j{+}1]$ for all $0 \leq j < n{-}1$.

\textbf{Case 1: $j \neq \mathit{lo}$ and $j+1 \neq \mathit{lo}$.}
$a_1[j] = a_0[j] \geq a_0[j{+}1] = a_1[j{+}1]$ by $\mathit{rev\_sorted}(a_0, n)$.

\textbf{Case 2: $j+1 = \mathit{lo}$ (the position after $\mathit{lo}$ receives more).}
Need $a_1[j] \geq a_1[\mathit{lo}] = x + 1$.
Since $j = \mathit{lo} - 1 < \mathit{lo}$, by Lemma~\ref{lem:search-correct}: $a_0[j] > x$.
Since $a_0[j]$ is an integer, $a_0[j] \geq x + 1$.
And $a_1[j] = a_0[j]$ (unchanged), so $a_1[j] \geq x + 1 = a_1[\mathit{lo}]$.

If $\mathit{lo} = 0$, this case does not arise (no position $j = -1$).

\textbf{Case 3: $j = \mathit{lo}$ (the position before $\mathit{lo}+1$).}
Need $a_1[\mathit{lo}] = x + 1 \geq a_1[\mathit{lo}{+}1]$.

If $\mathit{lo} = n-1$, this case does not arise.

If $\mathit{lo} < n{-}1$: $a_1[\mathit{lo}{+}1] = a_0[\mathit{lo}{+}1]$.
Since $\mathit{lo}{+}1 > \mathit{lo}$ and the array was reverse-sorted:
$a_0[\mathit{lo}{+}1] \leq a_0[\mathit{lo}] = x$.
Therefore $a_1[\mathit{lo}{+}1] = a_0[\mathit{lo}{+}1] \leq x < x + 1 = a_1[\mathit{lo}]$.

In all cases, $a_1[j] \geq a_1[j{+}1]$, so $\mathit{rev\_sorted}(a_1, n)$.
\end{proof}

\subsection{Task 4.3: Array Contents}

\begin{theorem}[Contents consistency]\label{thm:update-contents}
After \texttt{update}, the multiset of array values satisfies:
\[
  \mathit{contents}(a_1) = \mathit{contents}(a_0) - \mset{x} + \mset{x+1}
\]
where $x = a_0[i]$. This is equivalent to incrementing the value at position $i$
and re-sorting.
\end{theorem}

\begin{proof}
The function modifies only position $\mathit{lo}$, changing it from $x$ to $x + 1$.
\begin{itemize}
  \item $a_0[\mathit{lo}] = x$ (by Lemma~\ref{lem:search-correct}).
  \item $a_1[\mathit{lo}] = x + 1$ (the assignment).
  \item For all $j \neq \mathit{lo}$: $a_1[j] = a_0[j]$.
\end{itemize}
Therefore:
\begin{align*}
  \mathit{contents}(a_1)
  &= \mset{a_1[0], \ldots, a_1[n{-}1]} \\
  &= \bigl(\mset{a_0[0], \ldots, a_0[n{-}1]} - \mset{a_0[\mathit{lo}]}\bigr) + \mset{a_1[\mathit{lo}]} \\
  &= \mathit{contents}(a_0) - \mset{x} + \mset{x+1}.
\end{align*}

\textbf{Equivalence to increment-and-sort.}
Let $b$ be the array obtained by setting $b[i] \gets a_0[i] + 1$ and $b[j] = a_0[j]$
for $j \neq i$, then sorting $b$ in reverse order.
Since $a_0[\mathit{lo}] = a_0[i] = x$, the multisets are identical:
$\mathit{contents}(b_{\text{before sort}}) = \mathit{contents}(a_0) - \mset{x} + \mset{x+1} = \mathit{contents}(a_1)$.
Since sorting preserves multisets and $a_1$ is reverse-sorted (Theorem~\ref{thm:update-sorted}),
$a_1$ is exactly the reverse-sorted arrangement of $b$, i.e., $a_1 = \mathit{rsorted}(b)$.
\end{proof}

\subsection{Task 4.4: H-Index Correctness}

We first establish a key lemma about how the h-index changes when one value is incremented.

\begin{lemma}[H-index increment bound]\label{lem:hindex-incr}
Let $a'$ be obtained from $a$ by incrementing exactly one element by~1.
Then $\hindex(a', n) \in \{\hindex(a, n), \hindex(a, n) + 1\}$.
\end{lemma}
\begin{proof}
Let $h = \hindex(a, n)$ and $h' = \hindex(a', n)$.

\emph{$h' \geq h$:}
Suppose element $a[k]$ changes from $c$ to $c+1$.
For any threshold $v$:
\begin{itemize}
  \item If $v \leq c$: element $k$ satisfied $a[k] \geq v$ before and $a'[k] = c+1 \geq v$ after,
    so $\countge(a', n, v) = \countge(a, n, v)$.
  \item If $v = c+1$: element $k$ did not satisfy $a[k] \geq v$ before ($c < c+1$) but does after
    ($c+1 \geq c+1$), so $\countge(a', n, v) = \countge(a, n, v) + 1$.
  \item If $v > c+1$: element $k$ fails both before and after, so $\countge(a', n, v) = \countge(a, n, v)$.
\end{itemize}
In all cases, $\countge(a', n, v) \geq \countge(a, n, v)$.
Therefore $\countge(a', n, h) \geq \countge(a, n, h) \geq h$, so $h' \geq h$.

\emph{$h' \leq h + 1$:} We show $\countge(a', n, h{+}2) < h{+}2$.
Since $\countge(a', n, h{+}2) \leq \countge(a, n, h{+}2) + 1$ (at most one element increased
past $h{+}2$) and $\countge(a, n, h{+}2) \leq \countge(a, n, h{+}1) < h + 1$
(by monotonicity of $\countge$ and the h-index upper bound for $a$), we get
$\countge(a', n, h{+}2) \leq h + 1 < h + 2$.
Now suppose for contradiction that $h' \geq h+2$. By \ref{h:lower} for $a'$,
$\countge(a', n, h') \geq h'$, and by monotonicity $\countge(a', n, h{+}2) \geq
\countge(a', n, h') \geq h' \geq h{+}2$, contradicting the bound above.
Therefore $h' \leq h + 1$.
\end{proof}

\begin{theorem}[H-index correctness of \texttt{update}]\label{thm:update-hindex}
The return value of \texttt{update} equals $\hindex(a_1, n)$.
\end{theorem}

\begin{proof}
Let $h = \hindex(a_0, n)$ and $h' = \hindex(a_1, n)$.
By Lemma~\ref{lem:hindex-incr}, $h' \in \{h, h+1\}$.

By Corollary~\ref{cor:transition-equiv} applied to the reverse-sorted array $a_1$,
$h' = h + 1$ if and only if:
\begin{equation}\label{eq:incr-cond}
  \forall\, j < h{+}1.\; a_1[j] > j
  \qquad\text{and}\qquad
  \bigl(h{+}1 = n \;\vee\; a_1[h{+}1] \leq h{+}1\bigr).
\end{equation}

We analyze the two cases in the function's return logic.

\medskip\noindent
\textbf{Case A: $\mathit{lo} = h$ and $a_1[\mathit{lo}] = h + 1$.}

Since $\mathit{lo} = h$ and $a_1[h] = h + 1$:
\begin{itemize}
  \item Since $a_0[\mathit{lo}] = x$ and $a_1[\mathit{lo}] = x + 1 = h + 1$, we get $x = h$.
  \item For $j < h$: $a_1[j] = a_0[j]$ (unchanged, since $j \neq \mathit{lo} = h$).
    Since $\hindex(a_0, n) = h$ and $\mathit{rev\_sorted}(a_0, n)$,
    by Corollary~\ref{cor:transition-equiv}: $a_0[j] > j$ for all $j < h$.
    So $a_1[j] > j$ for all $j < h$.
  \item For $j = h$: $a_1[h] = h + 1 > h$, so $a_1[j] > j$.
  \item Combined: $\forall\, j \leq h.\; a_1[j] > j$, i.e., $\forall\, j < h{+}1.\; a_1[j] > j$.
  \item For the second part of~\eqref{eq:incr-cond}: If $h + 1 = n$, done.
    If $h + 1 < n$: $a_1[h{+}1] = a_0[h{+}1] \leq a_0[h] = x = h \leq h + 1$.
    (The first inequality is by reverse-sorting; the second by $x = h$.)
\end{itemize}
Both conditions of~\eqref{eq:incr-cond} hold, so $h' = h + 1$.
The function returns $h + 1 = h'$. \checkmark

\medskip\noindent
\textbf{Case B: $\mathit{lo} \neq h$ or $a_1[\mathit{lo}] \neq h + 1$.}

We show $h' = h$ by verifying the conditions of Corollary~\ref{cor:transition-equiv}
for $h$ on array $a_1$.

\textbf{First condition:} $\forall\, j < h.\; a_1[j] > j$.

For $j < h$ with $j \neq \mathit{lo}$: $a_1[j] = a_0[j] > j$ (by Corollary~\ref{cor:transition-equiv}
for $a_0$).

For $j = \mathit{lo}$ (if $\mathit{lo} < h$): $a_1[\mathit{lo}] = x + 1 > x \geq a_0[\mathit{lo}] = x$.
We need $a_1[\mathit{lo}] > \mathit{lo}$.
Since $\mathit{lo} < h$ and $a_0[\mathit{lo}] > \mathit{lo}$ (from h-index of $a_0$),
$a_1[\mathit{lo}] = a_0[\mathit{lo}] + 1 > \mathit{lo} + 1 > \mathit{lo}$.
Actually, since $a_0[\mathit{lo}] = x$ and $a_0[\mathit{lo}] > \mathit{lo}$,
we have $a_1[\mathit{lo}] = x + 1 > \mathit{lo} + 1 > \mathit{lo}$. \checkmark

\textbf{Second condition:} $h = n$ or $a_1[h] \leq h$.

If $h = n$: done.

If $h < n$ and $\mathit{lo} \neq h$: $a_1[h] = a_0[h] \leq h$ (from Corollary~\ref{cor:transition-equiv}
for $a_0$). \checkmark

If $h < n$ and $\mathit{lo} = h$: Then $a_1[h] = x + 1$, and by our case assumption $x + 1 \neq h + 1$,
so $x \neq h$. Since $a_0[h] = x$ and $a_0[h] \leq h$ (from Corollary~\ref{cor:transition-equiv}),
$x \leq h$ and $x \neq h$ gives $x \leq h - 1$, so $x + 1 \leq h$, i.e., $a_1[h] \leq h$. \checkmark

Both conditions hold, so $h' = h$ by Corollary~\ref{cor:transition-equiv}.
The function returns $h = h'$. \checkmark

\medskip
In both cases, the returned value equals $\hindex(a_1, n)$.
\end{proof}

\begin{remark}[Why the leftmost occurrence?]
The binary search finds the \emph{leftmost} position with value $x$ because incrementing
the leftmost position preserves the reverse-sorted order. If we incremented any other
position $p > \mathit{lo}$ with $a[p] = x$, then $a[p] = x + 1 > x = a[\mathit{lo}]$
would violate $a[\mathit{lo}] \geq a[p]$ for $\mathit{lo} < p$.
\end{remark}
